\section{A complete barrier-constrained nonlinear variational problem}
\label{sec:r11-completed-barrier-variational}

Retain the strict analytic box and background extension
(TC1)--(TC2), with the original arbitrary symmetric
middle-positive separately closed $(p,p)$-current $S$.
Set $q=n-p\geq1$, $d=q+1$, and fix an actual
$a\in\mathcal E^1_{S,\partial}(V)$.
Write $\widehat{\mathcal E}^1_a$ for the Borel relative
completion of Definition~\ref{def:r11-barrier-completion}.
The background remains $S$; only the potential class
is enlarged by the stated capacity completion.

The frozen predecessor is R4
\path{completed_variational_retraction.tex}.
It treats a bounded completed problem with a fixed
collar, and already distinguishes its constrained
variational condition from an arbitrary source equation.
The result here is an unbounded global first-energy
problem with one actual power-energy lower barrier,
using the new critical continuity and sharp midpoint
modulus.  It does not claim a new bounded Hilbert
projection theorem.

\subsection{Global geometry in the completed order interval}

For completed potentials with actual first-energy
lower barriers, define
\[
\begin{split}
 \widehat B(u_0,\ldots,u_q)
    &=\int_V(-u_0)\widehat N(u_1,\ldots,u_q),\\
 \widehat e(u)&=\widehat B(u^{[q+1]}),\qquad
 \widehat A(u)=-\widehat e(u)/d,\\
 \widehat I(u,v)
    &=\widehat e(u)+\widehat e(v)
       -\widehat B(u,v^{[q]})-\widehat B(v,u^{[q]}).
\end{split}
 \tag{BV1}
\]
The integrals are finite by (RC14).  They use the
actual Borel multipliers, with the joint pole set
annihilated by every coefficient.

\begin{lemma}[Completed energy geometry and separation]
\label{lem:r11-completed-first-energy-geometry}
The form $\widehat B$ is symmetric, nonnegative and
multilinear on the completed positive cone.
The function $\widehat A$ is concave along every
segment, and
\[
 \widehat A\left(\frac{u+v}{2}\right)
      -\frac{\widehat A(u)+\widehat A(v)}2
                    \geq c_q\widehat I(u,v)\geq0,
 \qquad c_q=\frac{1-2^{-q}}{2(q+1)}.
 \tag{BV2}
\]
Moreover
\[
 \widehat I(u,v)=0
       \quad\Longleftrightarrow\quad
 u=v\text{ locally outside a set of outer }S
                         \text{-capacity zero}.
 \tag{BV3}
\]
On $\widehat{\mathcal E}^1_a$ one has
$0\leq\widehat e(u)\leq e_1(a)$.
For a segment $u_t=(1-t)u+tv$,
\[
             \frac{d}{dt}\widehat A(u_t)
                    =\int_V(v-u)\,d\widehat N_q(u_t).
 \tag{BV4}
\]
All signed integrals in this formula are absolutely
convergent.
\end{lemma}

\noindent\textbf{Status: \texttt{DIRECT\_COROLLARY}}
of the new global critical completion theorem,
the actual first-energy geometry and its sharp
midpoint estimate.  The separation assertion is
proved explicitly to avoid assuming an ambient
classical representative.

\begin{proof}
Choose classical defining rows for all the
finitely many completed slots, retaining their
fixed actual barriers.  Formula (RC14) passes
each global mixed energy and every diagonal
energy to its completed value.  It therefore
passes symmetry, positive multilinearity,
the constant-one Holder inequality, concavity,
and (SM8).  For $u\geq a$, choose its defining
rows above $a$.  Actual energy monotonicity
gives $e_1(u_l)\leq e_1(a)$, and (RC14)
gives the stated bound.

If $\widehat I(u,v)=0$, the same convergence
gives $I(u_l,v_l)\to0$ for their classical
defining rows.  These rows have uniformly
bounded energy.  The all-local outer-capacity
estimate (EV8) makes every actual gap between
$u_l$ and $v_l$ tend to zero in capacity.
The defining convergence of each clipped row,
and the capacity-small deep sublevels of the
fixed barriers, give the same conclusion for
the completed gap.  The level-set triangle
inequality then makes each positive gap of
$u-v$ outer-null.  A countable union proves
the forward implication of (BV3).
Conversely, outer-null changes preserve every
completed current and every actual integrable
multiplier by (RC4)--(RC8), so the expression
in (BV1) is zero.

Finally positive multilinearity makes
$\widehat A(u_t)$ a finite polynomial in $t$.
Differentiate its $d$ symmetric slots.
This is (BV4), with integrability from the
completed mixed-energy bound.  No distributional
product of two unbounded differences is being
introduced by this scalar calculation.
\end{proof}

\begin{lemma}[Finite diffuse data integrate barrier-bounded capacity limits]
\label{lem:r11-barrier-diffuse-integration}
Let $\mu$ be a finite positive Radon measure on $V$
which vanishes on every Borel set of zero ordinary-open
outer $S$-capacity.  Suppose
\[
                              \int_V(-a)\,d\mu<\infty.
 \tag{BV5}
\]
If $u_j,u\in\widehat{\mathcal E}^1_a$ converge in the
defining local truncated outer capacities, then
\[
                           \int_V|u_j-u|\,d\mu\longrightarrow0.
 \tag{BV6}
\]
In particular this integral is independent of the
chosen relative representatives.
\end{lemma}

\noindent\textbf{Status: \texttt{NEW\_PROVED}}
in the unbounded barrier completion; the underlying
finite-measure capacity argument is detailed below.

\begin{proof}
A finite measure that vanishes on outer-null Borel
sets is sequentially absolutely continuous for this
outer capacity.  Otherwise choose Borel $E_j$ with
capacities tending to zero but $\mu(E_j)\geq\epsilon$,
and pass to a subsequence with summable capacities.
The limsup has outer capacity zero by countable
subadditivity.  On the other hand, continuity from
above of the finite measure applied to the decreasing
tail unions gives its mass at least $\epsilon$.
This is a contradiction.

The actual pole set of $a$ is outer-null and hence
$\mu$-null.  Off it the inequalities
$a\leq u_j,u\leq0$ imply
$|u_j-u|\leq-a$.
On a fixed compact set remove a deep sublevel of $a$;
its capacity and hence its $\mu$ mass can be made
small.  On the complement full values agree with
one fixed truncation, whose convergence in capacity
and the preceding sequential absolute continuity
give convergence in $\mu$ measure.
Exhaust $V$ by compact sets and use finiteness of
$\mu$ to obtain convergence in measure on all of $V$.
The integrable bound $-a$ in (BV5) then gives
$L^1(\mu)$ convergence: every subsequence has an
almost-everywhere convergent further subsequence,
and dominated convergence applies to it.
This proves (BV6) for the whole sequence.
The same null-set argument proves representative
independence.
\end{proof}

\begin{theorem}[A unique completed barrier-constrained maximizer]
\label{thm:r11-completed-barrier-maximizer}
For every datum satisfying Lemma
\ref{lem:r11-barrier-diffuse-integration}, the functional
\[
      \widehat{\mathcal F}_\mu(u)
            =\widehat A(u)-\int_Vu\,d\mu,
              \qquad u\in\widehat{\mathcal E}^1_a,
 \tag{BV7}
\]
has a maximizer, unique modulo local outer
$S$-capacity zero.  Its value lies between zero
and $\int(-a)\,d\mu$.
Denote its relative class by $\Pi_a(\mu)$.
Every classical maximizing sequence in the
actual interval $\{a\leq u\leq0\}$ is Cauchy
in all local outer capacities and converges to
this completed class.
The maximizer is characterized by the precise
variational inequality
\[
 \int_V(v-u)\,d\bigl(\widehat N_q(u)-\mu\bigr)
               \leq0
       \quad\bigl(v\in\widehat{\mathcal E}^1_a,\ u=\Pi_a(\mu)\bigr).
 \tag{BV8}
\]
The inequality is equivalent to maximality, and
all its integrals are absolutely convergent.
It is not asserted to be the unrestricted
current equation $\widehat N_q(u)=\mu$.
\end{theorem}

\noindent\textbf{Status: \texttt{NEW\_PROVED}}
for the unbounded completed order interval and
arbitrary original $S$.

\begin{proof}
The zero potential is admissible and has functional
value zero.  Since $\widehat e(u)\geq0$ and
$-u\leq-a$, every value in (BV7) is at most
$\int(-a)\,d\mu<\infty$.
For any completed $u$ its classical defining
rows converge in energy by (RC14) and in
$L^1(\mu)$ by (BV6).
Thus the supremum over the completed interval
equals the supremum over its actual classical
members.

Choose a classical sequence approaching this
finite supremum $L$.
Every term has first energy at most $e_1(a)$.
The midpoint remains in the same actual interval.
Its linear datum term cancels from the midpoint
gap, so (SM8) gives
\[
 c_q I(u_j,u_l)
       \leq L-\frac{\mathcal F_\mu(u_j)+
                              \mathcal F_\mu(u_l)}2\longrightarrow0.
 \tag{BV9}
\]
The local capacity estimate (EV8) makes the
sequence Cauchy in the defining topology.
The fixed-barrier completeness theorem gives
a limit $u\in\widehat{\mathcal E}^1_a$.
Global critical energy continuity (RC14) and
(BV6) give
$\widehat{\mathcal F}_\mu(u)=L$.
This proves existence without an ambient upper
regularization or a classical compactness claim.

If $u,v$ are both maximizers, (BV2) and the
affinity of the datum integral give
$c_q\widehat I(u,v)\leq0$.
By (BV3) they are the same relative class.
The same Cauchy argument and uniqueness apply
to every classical maximizing sequence.

For any completed competitor $v$, the segment
$u_t=(1-t)u+tv$ stays in the interval.
The derivative at $t=0$ is the left side of
(BV8), by (BV4) and (BV6).
Maximality gives its nonpositive sign.
Conversely concavity gives the upper tangent
inequality for (BV7); (BV8) therefore makes
every competitor's value at most the value
at $u$.  The data integrals are controlled
by (BV5), and all curvature integrals by
the completed mixed-energy bound.
This proves both the characterization and
the claimed integrability.
\end{proof}

\begin{theorem}[Data stability and retraction on the finite-mass image]
\label{thm:r11-completed-barrier-data-stability}
Let $\mu,\nu$ satisfy (BV5) and the same
finite BT-diffuse hypotheses, and define
\[
                    \delta_a(\mu,\nu)
                         =\int_V(-a)\,d|\mu-\nu|.
 \tag{BV10}
\]
Then
\[
 \widehat I(\Pi_a(\mu),\Pi_a(\nu))
       \leq\delta_a(\mu,\nu).
 \tag{BV11}
\]
In particular, if $\delta_a(\mu_j,\mu)\to0$,
the completed maximizers converge in every
local outer capacity, every original mixed
current, and all the fixed-barrier weighted
limits of (RC8) and (RC14).

If $b\in\widehat{\mathcal E}^1_a$ and its top
measure $\mu_b=\widehat N_q(b)$ has finite
total mass, then it satisfies the stated
datum hypotheses and
\[
                                \Pi_a(\mu_b)=b
 \tag{BV12}
\]
in the relative capacity quotient.
If the actual barrier has finite total top mass
and positive first energy, the constant one
in (BV11) is optimal for that same barrier,
domain and background.
\end{theorem}

\noindent\textbf{Status: \texttt{NEW\_PROVED}}
for the quantitative completed variational map.
The term retraction means exactly (BV12).

\begin{proof}
Write $u=\Pi_a(\mu)$ and $w=\Pi_a(\nu)$.
Use $w$ as a competitor in (BV8) for $u$,
and $u$ as a competitor for $w$.
Adding the two resulting inequalities gives
\[
 0\leq\widehat I(u,w)
       \leq\int_V(u-w)\,d(\nu-\mu)
       \leq\int_V(-a)\,d|\nu-\mu|.
\]
The last inequality uses $|u-w|\leq-a$.
This proves (BV11).  The two-competitor argument
is the frozen bounded retraction mechanism;
the newly justified unbounded integrals and
completed interval make it applicable here.

For clarity, the passage from $\widehat I\to0$
to local capacity convergence is also valid
for completed rows.  Approximate each pair
by classical rows with their fixed barrier.
Their energies and discrepancies converge
by (RC14).  Apply (EV8) to the classical
rows, use the level-set triangle inequality
and the defining capacity convergence, and
then decrease its intermediate gap threshold
to the requested one.  The resulting bound
is the same fixed-interval capacity modulus.
Thus (BV11) gives the stated convergence,
and (RC8), (RC14) give all the current
and weighted conclusions.

For (BV12), outer-capacity control of the
completed current on each compact stage
shows that $\mu_b$ is BT diffuse.
Also
$\int(-a)\,d\mu_b=\widehat B(a,b^{[q]})<\infty$.
Its finite total mass is the printed
additional assumption.
The completed tangent inequality for
$\widehat A$ gives
\[
 \widehat A(v)-\widehat A(b)
                    \leq\int(v-b)\,d\mu_b .
\]
Thus $b$ maximizes (BV7) for datum $\mu_b$.
Uniqueness proves (BV12).
Finally, if $\mu_a$ has finite total mass and
$e_1(a)>0$, take $\mu=\mu_a$ and $\nu=0$.
The retraction identity gives $\Pi_a(\mu_a)=a$,
and $\Pi_a(0)=0$ because the objective is
$-\widehat e/d\leq0$.
Then
$\widehat I(a,0)=e_1(a)=\delta_a(\mu_a,0)$,
proving the stated optimality.
\end{proof}

\begin{proposition}[A completed variational maximizer need not solve the source equation]
\label{prop:r11-completed-maximizer-source-obstruction}
For the precise smooth moment-Laplacian background
(DO6), with $n\geq2$ and $q=1$, take the
nonzero compact smooth source (DO9).
For every actual
$a\in\mathcal E^1_{S,\partial}(V)$, the
maximizer $\Pi_a(\mu)$ exists, but
\[
                            D(\Pi_a(\mu)\,S)\ne\mu .
 \tag{BV13}
\]
More generally no member of
$\widehat{\mathcal E}^1_a$ has that source.
Thus the implication from completed variational
maximality to an unrestricted original source
equation is \texttt{FALSE}, even after the
classical representative requirement is relaxed
to this relative completion.
\end{proposition}

\noindent\textbf{Status: \texttt{NEW\_PROVED}}
as the extension of the frozen compact-source
obstruction to the new barrier completion.
It refines the earlier actual-class obstruction
without replacing its proof.

\begin{proof}
The datum (DO9) is finite and dominated by a
fixed multiple of the full normalized BT
capacity, so it is BT diffuse.
Equation (DO4), with $s=1$ and $v=a$,
gives $\int(-a)\,d\mu<\infty$.
The existence theorem therefore applies.

Suppose a completed member $u$ had
$D(uS)=\mu$, and choose classical defining
rows $a\leq u_j\leq0$.
Their Gauss traces $g_j=u_j\circ\sigma$
are finite convex functions on the real
box $Q$, by the frozen trace theorem.
The actual barrier has a finite convex
Gauss trace on all of $Q$, by the
finite-trace argument in the proof of
Theorem~\ref{thm:r11-power-compact-source-obstruction}.
Hence on every compact subbox the $g_j$
have uniform upper and lower bounds.
Elementary convex slope bounds give
equicontinuity on each smaller compact
subbox.  A diagonal finite-net argument
therefore supplies a subsequence converging
locally uniformly to a finite convex $g$.

The completed weighted current convergence
in degree zero gives $u_jS\to uS$ on
every original compact test.
For each real $\phi\in C_c^\infty(Q)$,
the exact original calibration now yields
\[
 \langle t_*D(uS),\phi\rangle
     =\lim_j\int_Q g_j\Delta\phi\,dy
     =\int_Q g\Delta\phi\,dy
     =\langle\Delta g,\phi\rangle.
 \tag{BV14}
\]
The purported source has compact tropical
support.  The frozen R10 convex compact-source
rigidity makes $g$ affine and $\Delta g=0$.
Equation (BV14) would therefore give
$t_*\mu=0$, contradicting the nonzero
positive source (DO9).
This proves (BV13) throughout the completed
interval.
It uses the full-original degree-zero
convergence before testing the pushforward;
no toric-only definition of the currents
or auxiliary-field quotient is substituted.
\end{proof}

\paragraph{The remaining equation gate is concrete.}
The theorem proves existence, uniqueness, source
stability and a retraction for a nonlinear
constrained relative problem.
Its Euler condition is exactly (BV8).
Promoting that inequality to equality against
every original test requires additional
admissible variations or a different potential
class.  Proposition
\ref{prop:r11-completed-maximizer-source-obstruction}
shows that this promotion is not a technical
consequence of the completion already built.
The finite signed diffuse relative linear
Dirichlet problem uses the larger specified
energy class and remains compatible with
this obstruction.
